MHC-I presentation can vary among tissues even when the presenting allele and source protein are fixed. Existing predictors largely estimate tissue-blind peptide--MHC compatibility and therefore do not directly address this tissue-associated preference. We formulate a matched-pair ranking problem: given two peptides from the same parent protein under the same presenting restriction, determine which has stronger presentation evidence in a specified tissue. To remove a simple cross-tissue-frequency cue, the two peptides in every pair are additionally matched on their total number of positive tissue occurrences. The resulting benchmark contains 77 Human tissue--HLA tasks and 11 Mouse tissue--H-2 tasks, where H-2 is the mouse major histocompatibility complex. Both species use one fixed internal train/test split with 50 test pairs per task and the same three final-training seeds (20260704--20260706). Species-specific TissuePMHC configurations are selected by three-fold cross-validation confined to the training files before the fixed tests are opened. The occurrence audit finds zero mismatched pairs and identical positive-occurrence distributions between labels in both species and partitions. Tuned TissuePMHC reaches mean AUROC/AUPRC values of 0.8136/0.8126 on Human and 0.8194/0.8279 on Mouse. Tissue-independent external predictors remain near random in Human (AUROC 0.4874--0.5044) and Mouse (0.4917--0.5009). The tuning folds are not treated as strict-generalization estimates, and no peptide-separated evaluation is reported. These results assess tissue-conditioned peptide ranking for previously represented tissue--MHC tasks; they do not establish a causal tissue-processing mechanism, unseen-context generalization, independent-cohort validity, or clinical utility.
